% ===== §2 The Justice Typology and the Descent to Ontology =====
\section{The Justice Typology and the Descent to Ontology}
\label{p18:sec:typology}

\noindent
The phenomenon of \xref{p18:sec:phenomenon} is invisible to the standing theories of justice, and the section argues that the invisibility is systematic rather than incidental---a horizon of the field as a whole and not a gap in any one theory. The contemporary theory of justice is organised, this section contends, around an unstated ontology of the subject as a \emph{bearer}: a locus that has holdings, occupies a status, commands intelligibility, and can be formed. Every injustice the field knows is an injury to something the subject bears. The wrong this paper isolates is an injury not to anything borne but to the bearing---to the subject's activity of being a subject---and it is invisible precisely because the field's implicit ontology has nowhere to put an injury there. The claim, then, is not modest. It is that the theory of justice has a floor, that the floor is set one storey too high, and that a class of real wrongs has been falling through the gap the misplaced floor conceals. The section establishes this against the two most developed maps of justice's depth---Fraser's and Honneth's---and then installs the missing layer and defends the metaphysics that licenses it.

\subsection{Against the reigning maps of depth}
\label{p18:sub:typ-against}

The most serious contemporary attempts to say how deep injustice goes are Nancy Fraser's and Axel Honneth's, and the paper's claim is sharpened by saying exactly where it parts from each.

Fraser argues that justice has two irreducible dimensions, redistribution and recognition, to which she later adds a third, political dimension of representation; her ``perspectival dualism'' treats distribution and recognition as two analytical perspectives on any social arrangement, neither reducible to the other \parencite{fraserhonneth2003}. Honneth presses the opposite, monistic thesis: recognition is the single fundamental category, and distribution a derivative of it, so that a sufficiently differentiated theory of recognition can ground justice entire \parencite{honneth1995}. They divide over whether justice has one ground or two; they agree, decisively for our purposes, on the \emph{kind} of thing injustice injures. For both, the injured object is the subject's standing---what she is taken to be, how she is regarded, the social bases of her self-relation. Fraser's misrecognised subject and Honneth's disrespected subject are alike subjects whose \emph{regard} has failed; the quarrel is only over whether failed regard reduces to, or stands beside, failed distribution.

This is the agreement the present paper breaks. Both maps, however deep they reach, terminate at the subject's \emph{standing}---and standing, however damaged, presupposes a subject who stands. Neither map has a category for an injury to the subject's \emph{continued existence as a generative agent}, because both take that existence for granted as the substrate their injuries operate upon. Honneth comes nearest, and the nearness is instructive. He makes \emph{love}---the recognition proper to primary relationships---the structural core of ethical life, the sphere in which basic self-confidence is formed \parencite{honneth1995}. One might think that a theory placing love at the foundation would catch a harm done by love. It does not, and cannot, for a reason its own structure dictates: Honneth's love is recognition that \emph{builds}---it confers the self-confidence on which all later autonomy rests---and his pathology is always recognition \emph{withheld} or \emph{distorted}. He has no place for a love that, fully and undistortedly given, effaces by the very completeness of its giving. The harm this paper names is committed in the presence of perfect Honnethian recognition: the effaced party is loved, esteemed, accorded her rights, confident in herself---and effaced regardless.

\begin{claim}[The shared horizon of the recognition theorists]
Fraser's dualism and Honneth's monism, for all their opposition, share a horizon: both locate injustice, at its furthest reach, at the level of the subject's \emph{standing}---her distribution and her recognition---and both presuppose, beneath that level, a subject whose existence is not itself at issue. Neither can name a wrong done to the subject's continued existence as a generative agent, and Honneth's recognition-in-love, which might seem the exception, is precisely the recognition in whose full presence the wrong of effacement is still done. The recognition theorists mapped the standing of the subject to its floor; they did not reach the being beneath the standing.
\end{claim}

\subsection{The stratification, and where it stops}
\label{p18:sub:typ-strata}

Set the established categories out as what they are: not a list but a stratification, ordered by the depth at which each injures the subject, each layer reaching past the one above toward the subject herself.

\emph{Distributive} injustice injures \emph{holdings}---divisible goods, allocable and reallocable; its remedy is transfer, and what it cannot see is any wrong surviving a correct distribution. Because it is shallowest it is also imperialist: wrongs at every deeper layer are perennially misdescribed in its vocabulary of shares, which is why effacement, too, will tempt redescription as a maldistribution (a temptation \xref{p18:sec:definition} must resist).

\emph{Recognitive} injustice injures \emph{standing}---a good not distributable, since respect cannot be transferred as a holding can. Here is the first descent: a wrong that correct distribution leaves untouched, the wrong of the misrecognised who may be materially provided for and degraded all the same. Rawls already marks the pressure toward this layer from within the distributive paradigm itself, when he names the social bases of self-respect ``perhaps the most important primary good'' \parencite{rawls1971}---a holding that is not quite a holding, a distributive good whose content is recognition, the distributive paradigm straining at its own edge.

\emph{Epistemic} injustice injures \emph{intelligibility}---the subject in her capacity as a contributor to shared sense, her testimony discounted, her interpretive resources starved. This descends again: one may be well-provided and well-regarded and still be unable to make one's sense count.

\emph{Constitutive} injustice injures the \emph{conditions of formation}---the developmental and relational conditions under which a subject comes to be a subject at all. This is the furthest layer the field has reached, and the first that is genuinely ontological: what it injures is not anything the subject bears but her coming-to-bear.

\begin{claim}[The field terminates at formation]
The established categories form a single descent---holdings, standing, intelligibility, the conditions of formation---each detecting a wrong that survives the repair of the layer above it. But the descent halts at \emph{formation}: the furthest wrong the field can name is done at the threshold of subjecthood, to one not yet a subject. There is no category for a wrong done, beneath formation, to the \emph{continued being} of a subject already formed. The very monotonicity of the descent is the evidence that the floor is incomplete, for the series points past the storey on which the field has stopped.
\end{claim}

\subsection{The descent to ontology, and its price}
\label{p18:sub:typ-descent}

The missing layer is reached by one thesis: that for a relational subject, \emph{existence is an activity and not a state}. The field's implicit ontology pictures the subject as a substance that endures and bears properties, on which picture the gravest injury available is to prevent the substance from forming---hence the floor at constitutive injustice. The subject of intimate life, on this series' account, is not such a substance. Its relational being \emph{is} an activity: the continual self-constitution by which it generates itself as a generative subject in coupling with another. It does not first exist and then generate; its generating is the form its existence takes.

This thesis is the paper's largest commitment, and intellectual honesty requires stating its price rather than smuggling it in behind a metaphor. The objection that must be met is direct: \emph{is the identity of being with generation an argument, or merely a stipulation congenial to the conclusion?} A critic will say that ``to be is to generate'' is a definition chosen because it yields the wanted result, and that one could as easily hold the commonsense view---that the effaced subject plainly continues to exist (she breathes, she persists, she is the same person), and has merely been made unhappy or idle. On the commonsense view there is no ontological harm, only a distributive or eudaimonic one, and the whole edifice falls.

The reply is not to insist on the identity but to show what each side must pay for it. The commonsense view purchases its ease at a cost it does not advertise: it must hold that a subject's \emph{being} is independent of her \emph{activity}, that she would be the same being whether generating freely or reduced to the inert reception of another's generation---that, in the limit, a subject permanently stilled, permanently the mere site at which another's value is deposited, has lost nothing of her being and only something of her contentment. This the series denies, and denies for a reason given long before the present paper needed it: that a relational subject is \emph{constituted in the activity of relating}, so that the cessation of that activity is not a change in the subject's circumstances but a change in the subject. The identity of being with generation is not a free stipulation; it is the entailment of a relational ontology argued elsewhere on independent grounds, and its price---the rejection of the substance-subject who persists unchanged beneath all activity---is one the series has paid in full from its first paper. What the present section adds is only the observation that justice has not yet been made to reckon with that ontology's consequence: that if being is activity, then activity can be unjustly foreclosed, and the foreclosure is an injury to being.

\begin{claim}[The descent licensed by identity, not analogy]
Beneath constitutive injustice lies a final layer, ontological injustice, injuring a relational subject's \emph{being-as-generation}. The descent is licensed by an identity---for a relational subject, to be is to go on generating---and the identity is not a stipulation but the entailment of the relational ontology the series argues independently. Its price is the denial of the substance-subject who would persist unchanged beneath a foreclosed activity; a theory unwilling to pay that price keeps the commonsense subject and must redescribe effacement as mere discontent, at the cost of calling just a relation in which one party has been reduced to the inert site of another's generation. The layer is distinct in register from the constitutive: constitutive injustice wrongs the \emph{formation} of a not-yet-subject; ontological injustice wrongs the \emph{continued existence} of a subject already formed.
\end{claim}

\subsection{The typology, as a formal table}
\label{p18:sub:typ-table}

The descent can now be set out formally, and the table is itself an argument: its final column gives, for each layer, the reason it is irreducible to the layer above, so that read downward the table is a proof that the layers are distinct and the series cannot be collapsed into distribution.

\begin{center}
\renewcommand{\arraystretch}{1.5}
\footnotesize
\begin{tabularx}{\linewidth}{@{}>{\bfseries\color{rosedeep}\RaggedRight}p{1.75cm} >{\RaggedRight}p{2.4cm} >{\RaggedRight}p{2.5cm} X@{}}
\toprule
\textcolor{rosedeep}{Layer} & \textcolor{rosedeep}{Injury / locus} & \textcolor{rosedeep}{Detected by} & \textcolor{rosedeep}{Why not reducible to the layer above} \\
\midrule
Distributive
& holdings; \emph{what one has}
& maldistribution of a divisible good
& (base layer) \\

Recognitive
& status; \emph{what one is taken to be}
& misrecognition, denial of standing
& standing is not divisible; one may be justly provided for and degraded in regard (Rawls's ``social bases of self-respect'' already strains the distributive paradigm) \\

Epistemic
& intelligibility; \emph{whether one can be heard}
& testimonial / hermeneutical silencing
& sense-making is not a conferred status; one may be fully respected and still unable to make one's sense count \\

Constitutive
& conditions of formation; \emph{whether one can come to be}
& degradation of the conditions of subject-formation
& formation is prior to intelligibility; one may inhabit a rich epistemic order and still be barred, upstream, from forming as its subject \\

\rowcolor{pageblush}
\textcolor{rosedeep}{\textbf{Ontological}}
& \textcolor{rosedeep}{being-as-generation; \emph{whether one continues to generate}}
& \textcolor{rosedeep}{capture of the subject's own generative frequency (\xref{p18:sec:criterion})}
& \textcolor{rosedeep}{continuation is past formation; one may be formed, provided, regarded, and heard---and still stopped from continuing to generate as oneself} \\
\bottomrule
\end{tabularx}
\end{center}

\noindent
The locus column moves monotonically inward---has, is-taken-to-be, can-make-known, can-come-to-be, continues-to-be---and the inward movement is the descent made visible. The right-hand column shows that at each step the lower category detects a wrong that \emph{persists when the upper category's injury is repaired}: provision does not cure misrecognition; recognition does not cure silencing; audibility does not cure foreclosed formation; and---the step this paper adds---formation, with all the rest, does not cure foreclosed continuation. Each layer is the residue the layer above leaves uncured, and that residue is what makes the series irreducible.

\subsection{The irreducibility of the bottom layer}
\label{p18:sub:typ-irreducible}

The bottom row's irreducibility must be demonstrated and not merely tabulated, for it is the paper's claim to novelty. The demonstration has the form the table announces: in effacement, each of the four shallower criteria is \emph{satisfied-as-absent}, and the harm remains. Aggregate value rises and the effaced party holds no less---\emph{no distributive wrong}. She is regarded, esteemed, accorded her rights, confident in herself---\emph{no recognitive wrong, in any of Honneth's three spheres}. She is heard with full attention, her contributions credited---\emph{no epistemic wrong}. She is a mature subject whose formation was long complete---\emph{no constitutive wrong}. All four diagnostic conditions are absent; the harm is present, real, and grave. A wrong surviving the joint absence of all four cannot, as a matter of logic, be a compound of them.

\begin{claim}[Irreducibility]
Generative effacement is irreducible to the distributive, recognitive, epistemic, and constitutive injustices above it. The proof is that it is realised where all four shallower criteria are satisfied-as-absent---rising holdings, full recognition across all three Honnethian spheres, full audibility, completed formation---and the harm persists; were it a compound of the four it could not survive the absence of all four. The theory of justice that stopped at the fourth layer would be forced either to pronounce such a case just or to fall silent before it. The ontological layer is what gives justice a language for the subject who is given everything and ceases, nonetheless, to be.
\end{claim}

\subsection{A stress test: the case of chosen withdrawal}
\label{p18:sub:typ-stress}

A typology earns its categories by surviving the cases that would dissolve them, and one case presses hardest here. Suppose a subject \emph{chooses} to withdraw from generation---steps back, defers, makes herself the support rather than the source, in the manner traditionally praised as devotion or self-sacrifice. Is this effacement? If it is, the theory seems to pathologise every chosen self-subordination, an overreach that would discredit it; if it is not, the theory seems to admit that consent dissolves the harm, surrendering the very claim, central to this paper, that effacement survives sincere endorsement (\xref{p18:sub:phen-sincere}). The case appears to force a dilemma fatal either way.

It does not, and seeing why sharpens the concept. The dilemma trades on an equivocation in ``chooses to withdraw.'' Distinguish two things a subject might choose. She might choose a \emph{particular exercise of her generativity that takes the form of support}---directing her generation, at her own frequency, toward sustaining a shared project she has helped constitute as hers; this is not effacement at all, but generation, since her own dynamics are running, merely running toward a chosen end. Or she might choose a \emph{standing withdrawal of her generativity as such}---ceasing to generate at her own frequency, becoming the inert site of another's. Only the second is effacement, and the criterion of \xref{p18:sec:criterion} discriminates them exactly: the first leaves the subject's own frequency intact (she is the generator of her support), the second extinguishes it (her motion is now the echo of another's rhythm). The traditional praise of ``self-sacrifice'' conflates the two, and the conflation is precisely the ideological cover under which effacement passes for virtue. The theory does not pathologise the first; it \emph{names} the second, and names also the equivocation that has let the second hide inside the honour due the first.

\begin{claim}[Chosen support is not effacement; chosen self-erasure is]
The objection from chosen withdrawal dissolves on a distinction the criterion makes exact. To direct one's generativity, at one's own frequency, toward a freely constituted supporting role is generation, not effacement: the subject's own rhythm runs. To cease generating at one's own frequency, becoming the site of another's generation, is effacement whether or not it is chosen---and that it can be chosen, and honoured as sacrifice, is not a refutation but a confirmation, for it is exactly the wrong this paper holds can be completed in sincere endorsement. The theory pathologises no chosen support; it exposes the equivocation by which self-erasure has borrowed the prestige of support.
\end{claim}

\noindent
With the floor installed, its metaphysics defended, the layer tabulated, the bottom row shown irreducible, and the hardest case turned, the wrong has its locus and its place in the order of injustices. What it does not yet have is a definition exact enough to border it from its neighbours, a criterion able to detect it, a justice answering it, and a ground making that justice owed. To the definition the next section turns.
